\begin{problem}[Monomorphisms and kernels]\label{prob:vi15-p11}
Prove that $\varphi$ is a monomorphism iff $\ker\varphi=0$, iff every stalk map $\varphi_x$ is injective.
\end{problem}

\ifdefined\FullProblemDossiers
\paragraph{Why this problem matters.}
Unlike epimorphisms, monomorphisms are also sectionwise injective for sheaves of abelian groups.

\paragraph{Useful ideas.}
Translate the sheaf statement to local representatives or stalks whenever possible.  Keep track of whether the construction is sectionwise or requires sheafification.

\paragraph{Guided hints.}
Use the kernel universal property and stalk equality.

\begin{solution}
If $\varphi$ is a monomorphism, the inclusion $i:\ker\varphi\to\mathcal F$ satisfies
\[
\varphi i=0=\varphi 0.
\]
Monomorphy gives $i=0$, hence $\ker\varphi=0$. Conversely, if $\ker\varphi=0$ and $a,b:\mathcal H\to\mathcal F$ satisfy $\varphi a=\varphi b$, then $\varphi(a-b)=0$. The kernel universal property forces $a-b$ to factor through the zero sheaf, hence $a=b$. Therefore $\varphi$ is mono iff $\ker\varphi=0$.

Using
\[
(\ker\varphi)_x=\ker\varphi_x,
\]
the kernel sheaf is zero iff all stalk kernels are zero. This is equivalent to injectivity of every $\varphi_x$.
\end{solution}

\paragraph{What happened mathematically.}
Unlike epimorphisms, monomorphisms are also sectionwise injective for sheaves of abelian groups.

\paragraph{Extensions.}
Formulate the same statement for sheaves of modules.  Where meaningful, test the statement on a discrete space and then on the exponential-map example to see the difference between sectionwise and local behavior.

\paragraph{Related problems.}
Compare VI.15.P1 and VI.15.P2, the two preserved corpus problems.  VI/16 will reorganize these constructions into general exact-sequence language.

\paragraph{Provenance.}
New synthesis.
\else
\paragraph{Provenance.} New synthesis.
\fi
